\documentclass{article}
\usepackage{kotex}
\usepackage{graphicx}
\usepackage{booktabs}
\usepackage{amsmath}
\usepackage{multirow}
\usepackage{geometry}
\geometry{margin=1in}

\title{\textbf{\LARGE Controlling Contextual Embedding Bias via Noise Injection and Denoising: A Lightweight Approach}}
\author{NohGyubeom gnorm42@gmail.com}
\date{2025-03-24}

\begin{document}

\maketitle

\begin{abstract}
Pretrained contextual embeddings, such as BERT, often contain inherent biases reflecting social stereotypes \cite{sheng2019, kurita2019}. In this paper, we propose an effective and lightweight method to control and mitigate such biases by injecting noise into the embedding space and applying a denoising head. The proposed method is empirically evaluated on the Korean NSMC sentiment dataset \cite{nsmc2016}, and bias is measured using the Word Embedding Association Test (WEAT) \cite{caliskan2017, may2019} across gender, occupation, and nationality categories. Our approach not only mitigates the strength of biases but also alters their directional tendencies, presenting a practical solution for fairer NLP models. Furthermore, we discuss its applicability in resource-constrained environments, with potential integration into lightweight LLMs deployed in mobile or embedded systems \cite{sun2020mobilebert, lan2019albert}.
\end{abstract}

\section{Introduction}
Pretrained language models (PLMs) have become central to NLP tasks \cite{devlin2019, raffel2020}. However, these models inherit biases present in their training corpora \cite{bender2021, gehman2020}. While previous studies have primarily focused on biases in static embeddings, explorations into contextual embeddings remain relatively limited \cite{bolukbasi2016, manzini2019}. To bridge this gap, we propose a simple yet effective technique: injecting Gaussian and diffusion-based noise into BERT embeddings and recovering them through a denoising head.

\section{Related Work}
Bias issues in word embeddings have been extensively studied, revealing social biases such as gender and racial stereotypes in static embeddings like Word2Vec and GloVe \cite{mikolov2013, pennington2014, bolukbasi2016, manzini2019}. Methods such as hard debiasing and projection removal have been proposed to mitigate these biases \cite{bolukbasi2016}. However, contextual embeddings like BERT and GPT pose additional challenges due to their dynamic nature \cite{kurita2019}. Recent research has attempted to address this via fine-tuning or data augmentation \cite{zmigrod2019, zhao2018}, but often requires considerable computational resources. Our study introduces a novel lightweight alternative for effective bias control through noise injection and denoising.

\section{Methodology}
\subsection{Baseline Model}
We adopt a BERT-base model fine-tuned on the NSMC sentiment dataset as our baseline \cite{devlin2019, liu2019}.

\subsection{Noise Injection}
We consider two types of noise:
\begin{itemize}
\item \textbf{Gaussian Noise}: We add Gaussian noise with mean 0 and small variance directly to the embedding layer. The standard deviation is empirically set to 0.1.
\item \textbf{Diffusion Noise}: Noise is gradually injected over multiple steps following a forward diffusion process \cite{ho2020, song2020}.
\end{itemize}

The diffusion noise is defined mathematically as:
\begin{equation}
\mathbf{h}_t = \mathbf{h}_{t-1} + \epsilon_t, \quad \epsilon_t \sim \mathcal{N}(0, \sigma_t^2), \quad \sigma_t = \frac{t}{T} \times \sigma_{\text{max}}
\end{equation}
where $T$ is the total number of steps, and $\sigma_{\text{max}}$ is a hyperparameter controlling the noise scale.

\subsection{Denoising Head}
For both noise types, we apply a feedforward network consisting of two ReLU activation layers as the denoising head, jointly trained with the sentiment classifier to restore meaningful embeddings. The structure of the denoising head is expressed as:
\begin{equation}
\mathbf{h}_{\text{denoise}} = \text{ReLU}(W \mathbf{h}_{\text{noisy}} + b)
\end{equation}
The simple architecture of the denoising head is designed with computational efficiency in mind, allowing application even in resource-constrained environments \cite{sun2020mobilebert, lan2019albert}.

\section{Experiments}
\subsection{Dataset}
We use the NSMC sentiment dataset consisting of 150,000 Korean movie reviews, labeled with binary sentiment classes (positive/negative) \cite{nsmc2016}.

\subsection{WEAT Evaluation Setup}
For the WEAT evaluation, we utilize the following target and attribute word groups categorized into gender, occupation, and nationality domains:

\begin{itemize}
\item \textbf{Gender-Emotion Test}:
\begin{itemize}
    \item Target A: 남자, 남성, 아빠, 형, 아들
    \item Target B: 여자, 여성, 엄마, 누나, 딸
    \item Attribute 1: 행복, 기쁨, 사랑, 즐거움, 희망
    \item Attribute 2: 슬픔, 증오, 화, 불안, 절망
\end{itemize}

\item \textbf{Job Bias Test}:
\begin{itemize}
    \item Target A: 의사, 엔지니어, 프로그래머, 경찰, 변호사
    \item Target B: 간호사, 비서, 청소부, 교사, 요리사
    \item Attribute 1: 남자, 남성, 아빠, 형, 아들
    \item Attribute 2: 여자, 여성, 엄마, 누나, 딸
\end{itemize}

\item \textbf{Country Group Test}:
\begin{itemize}
    \item Target A: 한국, 일본, 중국, 러시아, 브라질
    \item Target B: 미국, 프랑스, 독일, 이집트, 인도
    \item Attribute 1: 행복, 기쁨, 희망
    \item Attribute 2: 슬픔, 화, 절망
\end{itemize}
\end{itemize}



\subsection{Training Setup}
All models are trained for 3 epochs using the AdamW optimizer with a learning rate of 5e-5 \cite{loshchilov2018}. The batch size is set to 32, and training is conducted on a single NVIDIA P100 GPU to demonstrate applicability in low-resource settings. The parameter count for the denoising head and noise injection modules is minimal compared to the original BERT, making the approach highly compatible with quantization \cite{banner2018, krishnamoorthi2018} or knowledge distillation \cite{hinton2015, tinybert2020, sun2020mobilebert} for deployment in mobile environments.

\subsection{Sentiment Classification Accuracy}
\begin{table}[htbp]
\centering
\caption{Sentiment Classification Accuracy Comparison}
\begin{tabular}{lc}
\toprule
\textbf{Model} & \textbf{Accuracy} \\
\midrule
Baseline  & 85.7\% \\
Gaussian Noise  & 85.0\% \\
Diffusion Noise & 84.2\% \\
\bottomrule
\end{tabular}
\end{table}

\subsection{WEAT Effect Size}
\begin{table}[htbp]
\centering
\caption{WEAT Effect Size}
\begin{tabular}{lccc}
\toprule
\textbf{Model} & \textbf{Gender-Emotion} & \textbf{Occupation} & \textbf{Nationality} \\
\midrule
Baseline  & 0.1250 & 0.6926 & -1.2141 \\
Gaussian Noise  & 0.2317 & 1.1806 & -0.4829 \\
Diffusion Noise & -0.5830 & -0.2260 & -0.4337 \\
\bottomrule
\end{tabular}
\end{table}

\section{Discussion and Conclusion}
In this study, we proposed a lightweight approach to mitigate inherent biases in pretrained BERT embeddings by injecting noise into the embedding space and restoring it using a denoising head. Experimental results demonstrated that both Gaussian and diffusion noise effectively adjusted bias intensity based on WEAT metrics, with diffusion noise notably reversing the bias direction.

This can be interpreted as the noise injection blurring certain biased patterns within the embeddings, while the denoising head encourages the recovery of more generalized representations. Although there was a slight drop in sentiment classification accuracy due to noise injection, the level of decrease was deemed acceptable considering the bias mitigation benefits. These results suggest that combining noise injection with a denoising head offers a simple yet powerful bias control mechanism.

Additionally, the proposed method only introduces a small number of parameters to the original BERT structure, confirming its feasibility for deployment in mobile environments when combined with quantization or knowledge distillation techniques.

Future work includes applying the proposed method to various PLMs (e.g., BERT variants, GPT) across different languages and datasets to assess generalizability. Furthermore, we plan to enhance bias control performance by extending the denoising head to Transformer-based architectures or integrating with techniques such as adversarial training and counterfactual data augmentation. Exploring the synergy with preference-based optimization methods (PPO, DPO) to refine bias control via human preference-based reward signals is also on our agenda.

\appendix
\section{Hyperparameters and Environment}
\begin{table}[htbp]
\centering
\caption{Hyperparameter Settings}
\begin{tabular}{lc}
\toprule
\textbf{Parameter} & \textbf{Value} \\
\midrule
Batch Size & 32 \\
Learning Rate & 5e-5 \\
Epochs & 3 \\
Noise Type & Gaussian / Diffusion \\
Noise Std (Gaussian) & 0.1 \\
Denoising Head & 2-layer ReLU FFN \\
Hardware & NVIDIA P100 GPU \\
\bottomrule
\end{tabular}
\end{table}

\section{WEAT Visualization}
\begin{figure}[htbp]
\centering
\includegraphics[width=0.8\textwidth]{weat_graph.png}
\caption{WEAT Effect Size Visualization across categories.}
\label{fig:weat_graph}
\end{figure}


% Bibliography remains unchanged

\begin{thebibliography}{99}

% Bibitems

\bibitem{devlin2019}
Jacob Devlin, Ming-Wei Chang, Kenton Lee, Kristina Toutanova. BERT: Pre-training of Deep Bidirectional Transformers for Language Understanding. NAACL-HLT, 2019.

\bibitem{raffel2020}
Colin Raffel et al. Exploring the Limits of Transfer Learning with a Unified Text-to-Text Transformer. JMLR, 2020.

\bibitem{bender2021}
Emily Bender, Timnit Gebru, Angelina McMillan-Major, Shmargaret Shmitchell. On the Dangers of Stochastic Parrots: Can Language Models Be Too Big? FAccT, 2021.

\bibitem{gehman2020}
Samuel Gehman, Suchin Gururangan, Maarten Sap, Yejin Choi, Noah A. Smith. RealToxicityPrompts: Evaluating Neural Toxic Degeneration in Language Models. EMNLP, 2020.

\bibitem{bolukbasi2016}
Tolga Bolukbasi, Kai-Wei Chang, James Y. Zou, Venkatesh Saligrama, Adam Kalai. Man is to Computer Programmer as Woman is to Homemaker? Debiasing Word Embeddings. arXiv preprint arXiv:1607.06520, 2016.

\bibitem{manzini2019}
Thomas Manzini, Yao Chong Lim, Yulia Tsvetkov, Alan W. Black. Black is to Criminal as Caucasian is to Police: Detecting and Removing Multiclass Bias in Word Embeddings. NAACL, 2019.

\bibitem{sheng2019}
Emily Sheng, Kai-Wei Chang, Premkumar Natarajan, Nanyun Peng. The Woman Worked as a Babysitter: On Biases in Language Generation. EMNLP, 2019.

\bibitem{kurita2019}
Keita Kurita, Nidhi Vyas, Ayush Pareek, Alan W Black, Yulia Tsvetkov. Measuring Bias in Contextualized Word Representations. ACL, 2019.

\bibitem{zmigrod2019}
Ran Zmigrod, Ekaterina Vylomova, Ryan Cotterell, Malvina Nissim. Counterfactual Data Augmentation for Mitigating Gender Stereotypes in Languages with Rich Morphology. ACL, 2019.

\bibitem{zhao2018}
Jieyu Zhao et al. Gender Bias in Coreference Resolution: Evaluation and Debiasing Methods. NAACL, 2018.

\bibitem{liu2019}
Yinhan Liu et al. Fine-tune BERT for Text Classification. arXiv preprint arXiv:1905.05583, 2019.

\bibitem{nsmc2016}
Naver Sentiment Movie Corpus. NSMC Dataset. 2016. \url{https://github.com/e9t/nsmc}

\bibitem{ho2020}
Jonathan Ho, Ajay Jain, Pieter Abbeel. Denoising Diffusion Probabilistic Models. NeurIPS, 2020.

\bibitem{song2020}
Yang Song, Jascha Sohl-Dickstein, Diederik P. Kingma, Abhishek Kumar, Stefano Ermon, Ben Poole. Score-Based Generative Modeling through Stochastic Differential Equations. ICLR, 2021.

\bibitem{loshchilov2018}
Ilya Loshchilov, Frank Hutter. Decoupled Weight Decay Regularization. ICLR, 2018.

\bibitem{banner2018}
Ron Banner, Yaniv Nahshan, Itay Hubara, Elad Hoffer, Daniel Soudry. Post Training 4-bit Quantization of Convolutional Networks for Rapid-Deployment. NeurIPS, 2018.

\bibitem{krishnamoorthi2018}
Raghuraman Krishnamoorthi. Quantizing deep convolutional networks for efficient inference: A whitepaper. arXiv preprint arXiv:1806.08342, 2018.

\bibitem{hinton2015}
Geoffrey Hinton, Oriol Vinyals, Jeff Dean. Distilling the Knowledge in a Neural Network. arXiv preprint arXiv:1503.02531, 2015.

\bibitem{tinybert2020}
Xiaoqi Jiao et al. TinyBERT: Distilling BERT for Natural Language Understanding. EMNLP, 2020.

\bibitem{sun2020mobilebert}
Zhiqing Sun, Hongkun Yu, Xiaodan Song, Renjie Liu, Yiming Yang, Denny Zhou. MobileBERT: a Compact Task-Agnostic BERT for Resource-Limited Devices. ACL, 2020.

\bibitem{lan2019albert}
Zhenzhong Lan, Mingda Chen, Sebastian Goodman, Kevin Gimpel, Piyush Sharma, Radu Soricut. ALBERT: A Lite BERT for Self-supervised Learning of Language Representations. ICLR, 2019.

\bibitem{caliskan2017}
Aylin Caliskan, Joanna J. Bryson, Arvind Narayanan. Semantics derived automatically from language corpora contain human-like biases. Science, 356(6334):183--186, 2017.

\bibitem{may2019}
Chandler May, Alex Wang, Shikha Bordia, Samuel Bowman, Rachel Rudinger. On Measuring Social Biases in Sentence Encoders. NAACL, 2019.

\bibitem{mikolov2013}
Tomas Mikolov, Kai Chen, Greg Corrado, Jeffrey Dean. Efficient Estimation of Word Representations in Vector Space. ICLR, 2013.

\bibitem{pennington2014}
Jeffrey Pennington, Richard Socher, Christopher D. Manning. GloVe: Global Vectors for Word Representation. EMNLP, 2014.

\end{thebibliography}

\end{document}
